\chapter{声音与声学}\label{ch:g12:sound-acoustics}

拨动吉他的低音弦，房间里便充满了低音 A，
而传到你耳中的，只有被挤压和拉伸的空气。
本章沿着这一思路展开：声音的传播速度、频率如何决定音调、功率如何决定响度，
以及为什么同一个音符在不同乐器上听起来并不相同。

\section{声音是一种压力波}

\begin{definition}[声波]\label{def:g12:sound-acoustics:sound-wave}
\emph{声波}\index{声波}是一种机械压力波
（见\cref{ch:g12:mechanical-waves}）：振动的表面会产生向外传播的
\emph{密部}\index{密部}（压强略高）和
\emph{疏部}\index{疏部}
（压强略低）。空气微团\emph{沿}传播方向振动
（声波是\emph{纵波}\index{纵波}），
且传播需要介质：置于真空罩下的铃铛便不再发声
（见\cref{ch:g10:signals-and-waves}）。
\end{definition}

\begin{omfigure}
\begin{tikzpicture}[font=\small]
  \fill[gray!55] (-0.75,-0.4) rectangle (-0.35,0.4);
  \draw[gray!75, fill=gray!20] (-0.35,0.28) -- (0.0,0.5) -- (0.0,-0.5)
    -- (-0.35,-0.28) -- cycle;
  \foreach \i in {0,...,69} \draw[gray!80, thin]
    ({0.2+\i*0.1+0.12*sin(\i*12.857)},-0.45) -- +(0,0.9);
  \draw[dashed, thin] (1.6,0.5) -- (1.6,0.7) (4.4,0.5) -- (4.4,0.7);
  \draw[omThm, Stealth-Stealth] (1.6,0.7) -- (4.4,0.7)
    node[midway, above] {$\lambda$};
  \draw[-Stealth] (1.6,-0.55) -- (1.6,-0.95) node[below] {密部};
  \draw[-Stealth] (3.0,-0.55) -- (3.0,-0.6) node[below] {疏部};
  \draw[-Stealth, omDef, thick] (7.45,0) -- (8.15,0) node[right] {$v$};
\end{tikzpicture}

{\small 扬声器驱动产生纵波：空气微团聚集成
一个个密部，相邻密部相距一个波长 $\lambda = v/f$。}
\end{omfigure}

\begin{proposition}[声速]\label{prop:g12:sound-acoustics:speed}
\index{声速}%
声音在空气中的传播速度约为 \qty{340}{m/s}（温度为
\qty{20}{\degreeCelsius}），在水中约为 \qty{1500}{m/s}，在钢中
约为 \qty{5000}{m/s}：介质越硬，
声速越快。\admitted
\end{proposition}

\begin{remark}[暖空气中声速更快]\label{rem:g12:sound-acoustics:temperature}
这些是实测值；对它们的预测将在本科一年级
分册中讲解。声音在较暖的空气中传播更快（\qty{331}{m/s}，
\qty{0}{\degreeCelsius} 时）——这正是大厅升温时管乐器走音的原因。
\end{remark}

\begin{example}[声波的尺度]\label{ex:g12:sound-acoustics:wavelengths}
关系式 $\lambda = v/f$ 可描述任意声音的尺度：在空气中，
\qty{20}{Hz} 的隆隆声的波长为 $340/20 = \qty{17}{m}$，约一栋楼的高度；而
\qty{20}{kHz} 的嘶嘶声只需 \qty{1.7}{cm} 的空间；米级的声音
很容易绕过门和拐角。
\end{example}

\section{音调、响度与耳朵}

\begin{definition}[音调与响度]\label{def:g12:sound-acoustics:pitch-loudness}
\emph{音调}\index{音调}是对频率的感知：
频率 $f$ 越高，音调越高；标准音 A 的频率为 \qty{440}{Hz}，
频率 $f$ 加倍则音调升高一个八度。\emph{响度}\index{响度}是对振幅的感知：
压强振荡越强，响度越大；
放大改变的是响度，而非音调。
\end{definition}

\begin{remark}[人耳的听觉范围]\label{rem:g12:sound-acoustics:audible}
人耳的听觉范围约为 \qty{20}{Hz} 至 \qty{20}{kHz}，
（见\cref{ch:g10:signals-and-waves}）；低于此范围为次声波，高于此范围为
超声波——蝙蝠的叫声可达 \qty{50}{kHz}，医用探头则达兆赫兹。
钢琴只取用中间的音域：\qtyrange{27.5}{4186}{Hz}。
\end{remark}

\section{声强与分贝标度}

\begin{definition}[声强]\label{def:g12:sound-acoustics:intensity}
\emph{声强}\index{声强} $I$ 是指
垂直于波面的一平方米面积上穿过的声功率：
$I = P/S$，单位为 \unit{W/m^2}。能听到的最弱声音的声强约为
$I_0 = \qty{1.0e-12}{W/m^2}$；痛觉约在 \qty{1}{W/m^2} 附近出现。
\end{definition}

\begin{proposition}[平方反比衰减]\label{prop:g12:sound-acoustics:inverse-square}
\index{平方反比律}%
小声源以功率 $P$ 向各方向均匀辐射，
则在距离 $r$ 处，$I = P/(4\pi r^2) \propto 1/r^2$：距离加倍
声强变为原来的四分之一。
\end{proposition}

\begin{proof}
整个功率 $P$ 穿过面积为 $4\pi r^2$ 的球面；两者相除即可。
\end{proof}

\begin{definition}[声强级]\label{def:g12:sound-acoustics:level}
人耳的听觉跨越十二个数量级，因此声强通常以
对数形式表示。声强 $I$ 的\emph{声强级}\index{声强级}为
\[ L = 10 \log_{10}(I/I_0), \qquad I_0 = \qty{1.0e-12}{W/m^2}, \]
单位为\emph{分贝}\index{分贝}（\unit{dB}）：听阈
为 \qty{0}{dB}，痛觉约在 \qty{120}{dB} 附近。
\end{definition}

\begin{proposition}[分贝运算]\label{prop:g12:sound-acoustics:db-arithmetic}
将 $I$ 乘以 $k$，声强级增加 $10\log_{10}k$ 分贝。特别地
$\times 10$ 则增加 \qty{10}{dB}；$\times 2$ 则增加
$10\log_{10}2 \approx \qty{3.0}{dB}$；$N$ 个互不相干的相同
声源的\emph{声强}相加，而非声强级相加，结果为
$L_1 + 10\log_{10}N$；小声源距离加倍（$I$
变为原来的四分之一）约减少 \qty{6.0}{dB}。
\end{proposition}

\begin{proof}
$10\log_{10}(kI/I_0) = 10\log_{10}(I/I_0) + 10\log_{10}k$;
距离情形：见\cref{prop:g12:sound-acoustics:inverse-square}，$k = 1/4$。
\end{proof}

\begin{method}[分贝运算]\label{met:g12:sound-acoustics:db}
\begin{enumerate}
  \item $L = 10\log_{10}(I/I_0)$；反过来，$I = I_0 \times 10^{L/10}$。
  \item 切勿直接相加声强级：先化为声强，相加，再化回
——或每加倍增加 $+\qty{3}{dB}$，每增大十倍增加 $+\qty{10}{dB}$。
  \item 小声源距离每加倍，声强级减少 \qty{6}{dB}。
\end{enumerate}
\end{method}

\begin{omfigure}
\begin{tikzpicture}[font=\small]
  \fill[omThm!12] (4.25,-0.14) rectangle (7.35,0.14);
  \draw[-Stealth] (-0.3,0) -- (7.7,0);
  \foreach \x/\l in {0/0, 1/20, 2/40, 3/60, 4/80, 5/100, 6/120, 7/140}
    \draw (\x,0.07) -- (\x,-0.07) node[below] {\l};
  \node[below] at (7.55,-0.07) {\unit{dB}};
  \foreach \x/\y/\t in {0/0.4/听阈，1.5/0.95/耳语，
      3/0.4/交谈，4/0.95/{繁忙街道}, 5.25/0.4/音乐会，
      6/0.95/痛阈，7/0.4/{喷气机起飞}}
    \draw[omDef] (\x,0.14) -- (\x,\y) node[above] {\t};
  \draw[omThm, thick] (4.25,-0.14) -- (4.25,0.14);
\end{tikzpicture}

{\small 日常声级阶梯：每个 \qty{20}{dB} 的台阶，
对应声强相差百倍；超过 \qty{85}{dB}（红色区域）
的暴露会损伤听力。}
\end{omfigure}

\begin{remark}[听力损伤]\label{rem:g12:sound-acoustics:damage}
听力损失取决于剂量：\qty{85}{dB} 下可安全接触八小时，而
每增加 \qty{3}{dB}（即声强加倍），可安全接触的
时间减半；受损的毛细胞无法再生。
\end{remark}

\section{音色与谐音}

\begin{definition}[谐音]\label{def:g12:sound-acoustics:harmonics}
频率为 $f_1$ 的周期性声音——其
\emph{基音}\index{基音}——一般是
若干频率为 $f_n = n f_1$，$n = 1, 2, 3, \dots$ 的正弦波之和，其
\emph{谐音}\index{谐音}。仅含基音的声音
为\emph{纯音}\index{纯音}，如音叉所发之声。
\end{definition}

\begin{definition}[频谱与音色]\label{def:g12:sound-acoustics:timbre}
声音的\emph{频谱}\index{频谱}是其谐音的柱状图：
每个 $f_n$ 对应一根柱，柱高为该谐音的
强度。耳朵把 $f_1$ 感知为音调，把各谐音强度的\emph{配比}
感知为\emph{音色}\index{音色}——即声音的
“色彩”：演奏同一音符的乐器共享 $f_1$，
区别仅在于各自的频谱。
\end{definition}

\begin{omfigure}
\begin{tikzpicture}[font=\small]
  \node at (1.45,1.0) {长笛}; \node at (5.85,1.0) {小提琴};
  \draw[-Stealth] (-0.1,0) -- (3.1,0) node[below] {$t$};
  \draw[omDef, thick, smooth, samples=120, domain=0:2.9]
    plot (\x, {0.55*sin(248*\x)});
  \begin{scope}[xshift=4.4cm]
    \draw[-Stealth] (-0.1,0) -- (3.1,0) node[below] {$t$};
    \draw[omDef, thick, smooth, samples=160, domain=0:2.9]
      plot (\x, {0.30*sin(248*\x)+0.18*sin(496*\x)
        +0.12*sin(744*\x)+0.08*sin(992*\x)});
  \end{scope}
  \begin{scope}[yshift=-3.3cm]
    \draw[-Stealth] (-0.1,0) -- (3.1,0) node[below] {$f$ (\unit{Hz})};
    \foreach \k/\h in {1/1.15, 2/0.10, 3/0.05} \fill[omProp!80]
      ({0.35+0.55*(\k-1)-0.09},0) rectangle ({0.35+0.55*(\k-1)+0.09},\h);
    \foreach \k/\l in {1/220, 2/440, 3/660, 4/880, 5/1100}
      \node[below] at ({0.35+0.55*(\k-1)},0) {\scriptsize \l};
    \begin{scope}[xshift=4.4cm]
      \draw[-Stealth] (-0.1,0) -- (3.1,0) node[below] {$f$ (\unit{Hz})};
      \foreach \k/\h in {1/0.80, 2/0.55, 3/0.38, 4/0.26, 5/0.18}
        \fill[omProp!80] ({0.35+0.55*(\k-1)-0.09},0)
          rectangle ({0.35+0.55*(\k-1)+0.09},\h);
      \foreach \k/\l in {1/220, 2/440, 3/660, 4/880, 5/1100}
        \node[below] at ({0.35+0.55*(\k-1)},0) {\scriptsize \l};
    \end{scope}
  \end{scope}
\end{tikzpicture}

{\small 同一 A 音（\qty{220}{Hz}）：波形（上）、频谱
（下）——基音与音调相同；小提琴的谐音
阶梯即其音色。}
\end{omfigure}

\begin{example}[为何能听出来电者]\label{ex:g12:sound-acoustics:two-instruments}
长笛的 A 音接近纯音；小提琴拉奏同样的
\qty{220}{Hz} 音时，在 \qtylist{440;660;880}{Hz} 处叠加强烈的谐音。
音调与响度相同，却没有人会把它们混淆：人声也是
一种你早已熟悉的频谱。
\end{example}

\section{振动的弦}

\begin{definition}[驻波]\label{def:g12:sound-acoustics:standing-wave}
被限制在弦两固定端之间的波可以形成
\emph{驻波}\index{驻波}：原地振动，
静止的点——\emph{波节}\index{波节}——与
振幅最大的点——\emph{波腹}\index{波腹}——交替排列。
不再有波动传播，弦以固定的模式振动。
\end{definition}

\begin{proposition}[弦的振动模式]\label{prop:g12:sound-acoustics:string-modes}
长为 $L$、两端固定的弦上，传播速度为
$v$ 的波，弦只有在能恰好容纳整数 $n$
个半波长、且半波长位于两端的波节之间时，弦才能稳定振动：
\[ \lambda_n = \frac{2L}{n}, \quad f_n = n\,\frac{v}{2L} = n f_1, \quad n = 1, 2, 3, \dots \]
最低模式 $f_1 = v/(2L)$ 为基音；其余模式
恰好是它的谐音。\admitted
\end{proposition}

\begin{remark}[诚实的证明在哪里]\label{rem:g12:sound-acoustics:honest}
受限的波恰好选取这些模式，这一结论源于
波与其反射波的叠加，这一计算在
本科一年级分册中完成。几何图像很有说服力：两端都必须保持
静止，因此必须恰好容纳整数个半波长。
\end{remark}

\begin{omfigure}
\begin{tikzpicture}[font=\small]
  \foreach \n/\y/\l in {1/0/{$f_1 = \dfrac{v}{2L}$},
      2/-1.6/{$f_2 = 2f_1$}, 3/-3.2/{$f_3 = 3f_1$}} {
    \draw[dashed, omExo] (0,\y) -- (4.6,\y);
    \draw[omDef, thick, smooth, samples=80, domain=0:4.6]
      plot (\x, {\y+0.5*sin(\n*39.13*\x)});
    \draw[omDef!50, thick, smooth, samples=80, domain=0:4.6]
      plot (\x, {\y-0.5*sin(\n*39.13*\x)});
    \fill (0,\y) circle (2pt) (4.6,\y) circle (2pt);
    \node[left] at (-0.25,\y) {$n = \n$};
    \node[right] at (4.95,\y) {\l};
  }
  \fill[omThm] (2.3,-1.6) circle (1.8pt);
  \fill[omThm] (1.533,-3.2) circle (1.8pt) (3.067,-3.2) circle (1.8pt);
  \node[omThm, below] at (2.3,-1.72) {波节};
\end{tikzpicture}

{\small 两端固定的弦的前三个模式：$n$
个半波长容纳于 $L$ 中，故 $\lambda_n = 2L/n$ 且 $f_n = nf_1$。}
\end{omfigure}

\begin{example}[小提琴的 A 弦]\label{ex:g12:sound-acoustics:violin}
小提琴的 A 弦长 $L = \qty{32.8}{cm}$，
发出 $f_1 = \qty{440}{Hz}$ 的音，沿弦的波速为
$v = 2Lf_1 = 2 \times 0.328 \times 440 \approx \qty{289}{m/s}$ ---
由张力与质量决定，由琴轴调谐。按下手指
使 $L$ 缩短，所有 $f_n$ 同时升高：频谱整体上移。
\end{example}

\begin{remark}[乐器为何不同]\label{rem:g12:sound-acoustics:instruments}
弦的激发方式——拨、拉、敲——决定
各模式所占的比重，琴体又对某些谐音的放大程度
高于其他谐音；激发方式与琴体共同决定频谱，故相同的
$f_1$ 从吉他、小提琴与钢琴中发出，各自都有自身的音色。
\end{remark}

\section{练习}

\begin{exercise}[$\star$]\label{exo:g12:sound-acoustics:1}
在空气中（$v = \qty{340}{m/s}$），求下列声音的波长：(a) 一个
\qty{55}{Hz} 的低音；(b) 一个 \qty{1.0}{kHz} 的蜂鸣声；(c) 一个
\qty{15}{kHz} 的尖啸声。比较最大与最小的波长。
\end{exercise}

\begin{exercise}[$\star$]\label{exo:g12:sound-acoustics:2}
将下列声音分为次声波、可闻声或超声波：\qty{8}{Hz}；
\qty{100}{Hz}；\qty{15}{kHz}；\qty{40}{kHz}；\qty{3}{MHz}。在
可闻声中，哪一个音调更低？
\end{exercise}

\begin{exercise}[$\star$]\label{exo:g12:sound-acoustics:3}
一名工人在 \qty{1.0}{km} 外敲击一根铁轨，计算声音
在钢中与在空气中传播所需的时间。用一只耳朵紧贴铁轨的
听者会听到什么？
\end{exercise}

\begin{exercise}[$\star$]\label{exo:g12:sound-acoustics:4}
(a) 对应 $I = \qty{1.0e-5}{W/m^2}$ 的声强级是多少？(b)
\qty{60}{dB} 时的声强是多少？(c) 说出与每个声强级接近的一种日常声音。
\end{exercise}

\begin{exercise}[$\star$]\label{exo:g12:sound-acoustics:5}
小提琴 G 弦发出 $f_1 = \qty{196}{Hz}$ 的音：给出其后的三个
谐音的频率。为什么 \qty{12}{kHz} 的哨声虽在声源处很丰富，
到达耳中时却像纯音？
\end{exercise}

\begin{exercise}[$\star\star$]\label{exo:g12:sound-acoustics:6}
警报器向各方向均匀辐射功率 $P = \qty{0.50}{W}$。
求距其 \qty{5.0}{m} 处的声强与声强级。
\end{exercise}

\begin{exercise}[$\star\star$]\label{exo:g12:sound-acoustics:7}
一把小提琴在你的座位处产生 \qty{68}{dB}。两把时的声强级是多少？
四把呢？要达到 \qty{78}{dB}，需要多少把小提琴？
\end{exercise}

\begin{exercise}[$\star\star$]\label{exo:g12:sound-acoustics:8}
距一台小型扬声器 \qty{2.0}{m} 处的声强级为 \qty{95}{dB}。
预测 \qty{8.0}{m} 处的声强级，它是否低于 \qty{85}{dB} 的界限？
\end{exercise}

\begin{exercise}[$\star\star$]\label{exo:g12:sound-acoustics:9}
一个 \qty{440}{Hz} 的音从空气进入水中，求其
在每种介质中的波长。潜水者听到的音高相同：哪一个量
由声源决定，哪一个量随之改变？
\end{exercise}

\begin{exercise}[$\star\star$]\label{exo:g12:sound-acoustics:10}
某频谱在 \qtylist{130;260;390;520}{Hz} 处有峰，
各峰高度依次减小：基音为多少？另一件乐器
以相同响度演奏同一个音：其频谱中哪些相同、
哪些不同？
\end{exercise}

\begin{exercise}[$\star\star$]\label{exo:g12:sound-acoustics:11}
长为 \qty{65.0}{cm} 的吉他弦发出
$f_1 = \qty{110}{Hz}$ 的音。求波速，再求
八度（\qty{220}{Hz}）与五度（\qty{165}{Hz}）所对应的振动长度。
\end{exercise}

\begin{exercise}[$\star\star\star$]\label{exo:g12:sound-acoustics:12}
一台医用探头向软组织发出 \qty{5.0}{MHz} 的超声波
（$v = \qty{1540}{m/s}$）：求其波长。回声经
\qty{26}{\micro s} 和 \qty{58}{\micro s} 后返回，分别来自某器官的两壁，
求两壁的深度及该器官的厚度。
\end{exercise}

\begin{exercise}[$\star\star\star$]\label{exo:g12:sound-acoustics:13}
户外舞台在 \qty{3.0}{m} 处产生 \qty{110}{dB} 的声级。将其视为小
声源，声级降至 \qty{85}{dB} 时距离为多少？
给出实际距离与此结果不同的两个原因。
\end{exercise}

\begin{exercise}[$\star\star\star$]\label{exo:g12:sound-acoustics:14}
求痛阈（\qty{1}{W/m^2}）与 $I_0$ 之比。
取鼓膜面积为 \qty{0.5}{cm^2}，计算两个阈值下
耳朵接收到的功率，并评述耳朵作为探测器的特点。
\end{exercise}

\begin{exercise}[$\star\star\star$]\label{exo:g12:sound-acoustics:15}
单个相同的节日扬声器在调音台处的声级为 \qty{80}{dB}。
若开启 $2$ 台，声级为多少？$10$ 台呢？$100$ 台呢？
每多 \qty{+10}{dB} 需付出什么代价？这对响度竞赛说明了什么？
\end{exercise}

\section{问题：弦乐器制作师的工坊}

\begin{problem}[{周末问题——弦乐器制作师的工坊：吉他弦逐品分布，
其泛音调成音色，演奏会保持在危险线以下，
以及用超声波寻找
木材内部瑕疵}]
\label{pb:g12:sound-acoustics:1}
一位弦乐器制作师正在完成一把吉他。其 A 弦的振动长度
$L = \qty{65.0}{cm}$，须发出 $f_1 = \qty{110}{Hz}$ 的音。取
\qty{340}{m/s} 为空气中的声速，并在本阶段不加证明地承认 $f_1 = v/(2L)$
（\cref{prop:g12:sound-acoustics:string-modes}），其中 $v$ 为
弦上的波速。

\medskip
\noindent\textbf{第一部分——划分指板。}
\begin{enumerate}
  \item 求弦上的波速 $v$。
  \item 求基音在\emph{弦上}的波长，
        再求 \qty{110}{Hz} 声音\emph{在空气中
        }的波长。同一频率下，为何二者不同？
  \item 八度（\qty{220}{Hz}）对应的振动长度为多少？
        第十二品应位于何处？用 $L$ 的分数表示。
  \item 五度（\qty{165}{Hz}）对应的振动长度为多少？
  \item 每品使音高升高一个半音，倍数为
        $2^{1/12}$。用对数检验 \qty{110}{Hz} 与 \qty{220}{Hz} 之间
        相隔几个半音？第一品处的振动长度为多少？
\end{enumerate}

\medskip
\noindent\textbf{第二部分——泛音与音色。}
\begin{enumerate}[resume]
  \item 列出空弦 A 弦的前五个泛音。
  \item 一位小提琴手的 A 音为 \qty{440}{Hz}。它是吉他弦的第几次
        泛音，与基音相隔何种音程？
        \qty{330}{Hz} 呢？
  \item 制作师先在琴桥附近拨弦，再在指板上方拨弦：
        音高相同而音色不同，试用频谱加以解释。
  \item 拨弦时若在中点轻触琴弦，使该处成为波节。
        哪些泛音得以保留？听到的音高为何？
  \item 原则上，空弦的泛音中有多少个
        落在可听范围之内？
\end{enumerate}

\medskip
\noindent\textbf{第三部分——演奏会。}
在吉他首场音乐会上，每件乐器作为小声源，
声功率为 $P = \qty{6.0e-3}{W}$，
向各方向均匀辐射。
\begin{enumerate}[resume]
  \item 求距单件乐器 $r = \qty{3.0}{m}$ 处的声强。
  \item 由此推得该处的声级。
  \item 四重奏演出时，四件同类乐器在同一位置
        产生的声级为多少？为何不是原声级的四倍？
  \item 十二件乐器同时演奏：\qty{3.0}{m} 处的声级为多少？
  \item 听众须坐在离乐队多远处，合奏的声级
        才会降至 \qty{85}{dB}？
  \item 安全规则：\qty{85}{dB} 的声级可安全收听八小时，且
        每增加 \qty{+3}{dB}，安全收听时间减半。第 14 题中
        的前排听众最多可安全停留多久？
\end{enumerate}

\medskip
\noindent\textbf{第四部分——木材中的瑕疵。}
上漆之前，制作师用 \qty{2.0}{MHz} 的探测器检测
一块厚 \qty{45}{mm} 的琴颈坯料；该木材中取 $v = \qty{5.0e3}{m/s}$。
\begin{enumerate}[resume]
  \item 求超声波在木材中的波长。探伤用的波为何必须是超声波，
        而不能是可听声？
  \item 求来自坯料背面的回声延迟时间。
  \item 回声在 \qty{7.2}{\micro s} 后到达：瑕疵深度为多少？
  \item 写出工坊卡片：弦上的波速、第十二品与第一品对应的振动长度、
        安全收听距离与瑕疵深度。
\end{enumerate}
\end{problem}
